\documentclass[11pt]{article}
\usepackage[T1]{fontenc}
\usepackage{lmodern}
\usepackage[margin=2.2cm]{geometry}
\usepackage{amsmath,amssymb}
\usepackage{booktabs,tabularx,multirow,float,array}
\newcolumntype{Y}{>{\raggedright\arraybackslash}X}
\usepackage{enumitem}
\usepackage{xcolor}
\usepackage[hidelinks]{hyperref}
\usepackage{microtype}
\setlist{itemsep=2pt,topsep=4pt}
\usepackage{xurl}
\urlstyle{tt}
\newcommand{\file}[1]{\path{#1}}
\newcommand{\todo}[1]{\textcolor{red!70!black}{[#1]}}

\title{Synthetic Maize and Sorghum Datasets:\\Single-View Captures and Occluded Row Views}
\author{}
\date{October 2026}

\begin{document}
\maketitle

\begin{abstract}
We describe how four synthetic plant datasets were generated: single-view captures of isolated maize and sorghum plants,
and \emph{row views} in which the same plants stand in a crop row between two neighbours and are photographed by a ground
rover. Every plant is procedurally generated, so its full parameter record is known. The single-view sets pair each
rendered view (RGB, depth, back-projected point cloud) with the plant's parameters. The row sets pair an occluded
observation of a plant inside a row (RGBA, depth, point cloud of all visible plants) with the ground truth of that plant
alone, rendered from the same camera with the same leaf shapes. Neighbouring plants in a row may touch but never
intersect; leaves that would cross another plant are bent around it by a contact solver.
\end{abstract}

\tableofcontents

%======================================================================
\section{Overview}

\begin{table}[H]
\centering\footnotesize
\begin{tabularx}{\linewidth}{>{\raggedright\arraybackslash}p{2.3cm}YYYY}
\toprule
 & Maize single view & Sorghum single view & Maize row view & Sorghum row view \\
\midrule
Plants & 15{,}000 & 15{,}000 & 15{,}000 rows & 15{,}000 rows \\
Split (plants or rows) & 10{,}500 / 2{,}250 / 2{,}250 & 10{,}500 / 2{,}250 / 2{,}250 & same as single view & same as single view \\
Images per plant/row & 10 views & 10 views & 10 orientations & 10 orientations \\
Total images & 150{,}000 & 150{,}000 & 150{,}000 & 150{,}000 \\
Camera & per-plant fitted distance, random directions & 2.5\,m, Fibonacci directions & rover, 1.0\,m, level & rover, 2.5\,m, level \\
Input cloud & 8{,}192 pts & --- (labelled scan copy) & 24{,}576 pts & 24{,}576 pts \\
Target & parameters & parameters & main plant alone: RGBA, depth, 8{,}192-pt cloud, parameters & same \\
\bottomrule
\end{tabularx}
\caption{The four datasets.}
\end{table}

\noindent Locations (relative to \file{/media/microway/back_up_2/}):
\begin{itemize}
\item maize single view: \file{MeshPointCloudCapture/Maize/{train,val,test}/}
\item sorghum single view: \file{EvoEngine/out/build/x64-Release/EvoEngine_App/Sorghum_15K/Sorghum_15K_noobj/}
\item maize row view: \file{RoverRows/maize_row/{train,val,test}/}
\item sorghum row view: \file{RoverRowsSorghum/Sorghum_row/{train,val,test}/}
\end{itemize}

All images are $1024\times1024$ with a $40^\circ$ vertical field of view (so $f_x=f_y=512/\tan 20^\circ=1406.708$\,px,
$c_x=c_y=512$). Coordinates follow the OpenGL convention: $y$ is up and a camera looks down its $-z$ axis.

%======================================================================
\section{Procedural plants}

\subsection{Maize}
The 15{,}000 maize plants (\file{MeshPointCloudCapture/baseline_15000}, \file{plant_0000}\ldots\file{plant_14999}) were
made with the procedural maize generator \file{ProceduralModelMaize-newgen} through its batch ``Generate Plant Set''
export. Plant $i$ is generated with random seed $i$. Each plant is written as
\begin{itemize}
\item \file{plant_XXXX.obj}: the mesh, with one \texttt{usemtl} block for the stem, one per leaf in leaf-id order, then the
tassel (shared \file{maize.mtl});
\item \file{plant_XXXX.xml}: the full generator record (tassel, tiller, and per leaf: distance along the stem,
leafLength, leafWidth, leafAzimuthDeg, leafAngle, droopiness, spline points, width taper, twist, curl, waves, noise,
midrib, ligule, sheath).
\end{itemize}
Traits are drawn from Gaussians whose means and standard deviations are set in \file{maize_generator_config.xml}
(e.g.\ height $0.832\pm0.104$\,m, leaf count $12.35\pm2.57$, base leaf length $0.448\pm0.057$\,m, base leaf width
$0.066\pm0.010$\,m, base leaf angle $43.98^\circ\pm10.13^\circ$), coupled through four latent traits (vigour,
slenderness, posture, texture) with correlation strength $0.735$, with bottom/middle/top canopy curves and per-leaf
jitter (scale $0.687$). The resulting plants are $0.62$--$1.44$\,m tall (median $1.01$\,m).

\subsection{Sorghum}
The 15{,}000 sorghum plants (\file{/media/microway/back_up_1/Alloy/SorghumData}, \file{Sorghum_0}\ldots\file{Sorghum_14999})
were generated by EvoEngine's \file{SorghumDataGeneratorApp} with the descriptor
\file{SorghumGenerator/Random.sg}, in batches of 500 plants (\file{generate_batched.sh}); plant $i$ uses seed $i$.
\file{SorghumGenerator::Apply} draws the plant heading uniformly in $[0^\circ,360^\circ)$, the stem tilt from a Gaussian,
the leaf count (clamped to $[2,128]$) and the internode length, and sets the stem length to
$\text{internode}\times n_\text{leaves}/(1-s_0)$ with $s_0$ the first leaf's starting point. Per leaf $k$ at position
$t=k/(n-1)$ it evaluates curves over $t$ for starting point, length, width, roll (alternating sides:
$\text{roll}_k = (k \bmod 2)\cdot180^\circ + \text{roll}(t)$), branching angle, bending, curling and waviness. No panicle
is generated. Main descriptor values: stem tilt $0\pm1.5^\circ$, internode length $0.04\pm0.03$\,m, stem width
$0.014\pm0.005$\,m, leaf amount $20\pm2$. The generated plants have 4--35 leaves (mean 19.5) and stems of
0.19--1.74\,m. Per plant the generator writes
\begin{itemize}
\item \file{Sorghum_N.obj}: the mesh (one stem block and one block holding every leaf, front and back faces);
\item \file{Sorghum_N_nc.ply}: an occlusion-free labelled point cloud (each organ scanned on its own and merged; fields
\texttt{x y z type\_index instance\_index leaf\_index}, leaf\_index 0 = stem);
\item \file{Sorghum_N.ply}: a simulated gantry scan (with occlusion);
\item \file{Sorghum_N_spline.yml}: seed, stem length and direction, and per leaf starting point, length, roll angle,
branching angle, waviness frequency and period start, plus the centre/left/right spline point lists;
\file{Sorghum_N.yml}: the binary spline record.
\end{itemize}

%======================================================================
\section{Train / validation / test splits}

Both species are split \emph{by plant} into 10{,}500 / 2{,}250 / 2{,}250 plants (70/15/15\,\%), so that all views and
all rows of one main plant fall in one split. The held-out sets are biased toward unusual phenotypes.

\paragraph{Maize.} \file{split_maize.py} describes each plant by 19 scalar features of its XML (6 constant tassel
features are dropped), standardises them robustly ($z=(x-\text{median})/(1.4826\,\text{MAD})$) and scores each plant by
its Mahalanobis distance from the median, $\sqrt{z^\top(\Sigma+10^{-6}I)^{+}z}$. Plants are sorted from most to least
extreme and dealt alternately to test and validation until both hold 2{,}250; the rest is train (seed 0). The split is a
pure tail split: train scores $1.56$--$4.60$, validation $4.60$--$23.3$, test $4.60$--$32.2$. Lists:
\file{MeshPointCloudCapture/Maize/plants_{train,val,test}.txt}.

\paragraph{Sorghum.} The split is the one shipped with the sorghum single-view set,
\file{Sorghum_15K_noobj/assignment.csv} (columns \texttt{plant, split, extremeness, rank}), with the features in
\file{features.csv}. The extremeness reproduces (to the 4 decimals stored) as the Mahalanobis distance from the mean,
with sample covariance, over six features: leaf count, stem length, mean and maximum leaf length, and mean and standard
deviation of leaf roll; rank is its empirical percentile. Unlike maize, the hold-out is not a pure tail: the probability
of being held out rises steadily with rank (from 0.3\,\% in the least extreme 5\,\% of plants to 70\,\% in the most
extreme 5\,\%), and validation and test have the same distribution (median rank 0.75 each, train 0.37).
\todo{The script that wrote assignment.csv was not found; the selection rule above is inferred from the file.}
For the row set the lists were written from \file{assignment.csv}, sorted by plant id:
\file{RoverRowsSorghum/Sorghum_row/splits/plants_{train,val,test}.txt}.

%======================================================================
\section{Single-view captures}

\subsection{Common renderer}
Both species are rendered off-screen with \texttt{pyrender} (EGL), $1024\times1024$, $40^\circ$ vertical FOV, on a grey
background (0.5). The camera at $\mathbf{e}$ looks at the plant's bounding-box centre $\mathbf{c}$ with world up $(0,1,0)$.
Lighting: ambient $0.15$ and three white directional lights aimed at $\mathbf{c}$, from directions $(1,1,1)$ (intensity
4.0), $(-1,0.5,-0.5)$ (2.0) and $(0,1,0)$ (1.5).

\paragraph{Depth encoding.} With $z$ the linear depth of a pixel and near/far planes $n,f$, the stored value is
$v=\lfloor \operatorname{clip}((z-n)/(f-n),0,1)\cdot(2^{32}-1)\rfloor$, written as a 4-channel PNG with
$R,G,B,A$ = bits 31--24, 23--16, 15--8, 7--0 of $v$; $v=0$ is background. Decoding:
$z = n + \tfrac{v}{2^{32}-1}(f-n)$.

\paragraph{Back-projected point cloud.} Every pixel $(u,v)$ with depth $z>0$ gives the camera-frame point
\begin{equation}
\mathbf{p}_\text{cam} = \Big(\tfrac{u-c_x}{f_x}z,\; -\tfrac{v-c_y}{f_y}z,\; -z\Big),
\label{eq:backproject}
\end{equation}
and a fixed number of these points is drawn uniformly at random without replacement (all of them if fewer). Clouds are
binary PLY with double-precision $x,y,z$.

\subsection{Maize single view}
\file{capture_maize.py}. For each plant, 10 view directions are drawn uniformly on the sphere (normalised Gaussian
vectors, generator seeded with $(0,\text{plant})$). The camera distance is fitted to the plant: with $r$ half the
bounding-box diagonal, $\text{dist}=1.05\,r/\sin 20^\circ$, near $=\max(0.001,\text{dist}-1.5r)$, far $=\text{dist}+1.5r$.
A fourth directional light (intensity 1.5) sits at the camera. Materials come from \file{maize.mtl}: textured leaf, stem
and tassel (tassel tinted $(0.95,0.91,0.55)$), metallic 0, roughness 0.9, double-sided.

Per view (\file{Maize/<split>/plant_XXXX_VV/}): \file{rgb.png}; \file{rgba.png} (alpha 255 where the plant is, background
cleared); \file{depth.png}; \file{pointcloud_cam.ply} (8{,}192 points, Eq.~\ref{eq:backproject} at pixel centres, seed
$(0,\text{plant},\text{view})$); \file{camera_pose.json} (position, quaternion, camera-to-world and world-to-camera
matrices, view direction, intrinsics, near, far, depth encoding, number of points, plant centre, source files); and
\file{maize_XXXX_spline.xml}, the trimmed parameter file described next.

\paragraph{Maize parameter file.} The source XML is reduced (\file{trim_params.py}) to the plant height (highest mesh
point above the ground $y=0$), the leaf count, and per leaf: id, distance along the stem, leafLength, leafWidth,
leafAzimuthDeg, leafAngle and droopiness. \file{leaf_curvature.py} adds per-leaf curvature: the midrib B-spline is
rebuilt from leafLength, leafAngle, droopiness and the spline point count (cubic, clamped, target elevation
$\text{leafAngle}+\text{droopiness}\cdot t^2$) and integrated over the blade, arc-length $s\in[0.22,1]$, to give the bend
angle $\theta=\int|\kappa|\,ds$, the blade length $S$ (10-node Gauss--Legendre) and the mean curvature $\theta/S$; a
plant summary gives $\sum\theta/\sum S$, the mean per-leaf curvature and the mean bend angle.

\subsection{Sorghum single view}
\file{capture.py}. The 10 view directions are a Fibonacci sphere, $y_i=1-2(i+\tfrac12)/10$,
$\phi_i=i(\sqrt5-1)\pi+\epsilon_i$ with a small fixed jitter $\epsilon_i$, so every plant is seen from the same 10
directions (five above and five below the horizon). The camera stays 2.5\,m from the bounding-box centre; near 0.025,
far 50. One leaf texture (\file{leafSurface.png}) covers the whole plant, metallic 0, roughness 1.0; there is no
headlight.

Per view (\file{Sorghum_15K_noobj/<split>/Sorghum_N_VV/}): \file{rgb.png}, \file{depth.png}, \file{camera_pose.json}
(position, quaternion, camera-to-world and world-to-camera matrices), \file{Sorghum_N_nc.ply} (the labelled
occlusion-free cloud, world frame), \file{Sorghum_N_nc_cam.ply} (the same points in the camera frame; labels not kept),
\file{Sorghum_N_spline.yml} and \file{Sorghum_N.yml}. No mesh is included.

%======================================================================
\section{Row views}

A row view shows a \emph{main} plant standing in a crop row between two neighbours, as seen by a ground rover driving
along the row. The ground truth is the main plant alone, with exactly the leaf shapes it has in that row, rendered from
the same camera. Rows were generated by \file{RoverRows/tools} (\file{rover_run.py} drives \file{rover_row.py}, which
calls the contact solver \file{rover_solve.js}/\file{solver.js}).

\subsection{Row composition}
Rows are built within one split from its plant list $L$: row $i$ has main plant $L_i$, right neighbour $L_{i+1}$ and left
neighbour $L_{i+2}$ (indices modulo $|L|$), so there is one row per plant and every plant is a main plant once. The world
frame is the main plant's: its stem base at the origin, $y$ up, the row along $x$, the rover on the $+z$ side. The right
neighbour's stem base is at $(+d,0,0)$, the left one at $(-d,0,0)$, with spacing $d=0.15$\,m between stem bases.

\paragraph{Orientations.} Each row is rendered in 10 orientations: the main plant is turned about the vertical axis
through its stem base by $\theta_\text{main}\in\{0^\circ,10^\circ,\dots,90^\circ\}$, and each neighbour by
$\theta_\text{main}+\delta$ with $\delta\sim U(-15^\circ,15^\circ)$, drawn for the right then the left neighbour with a
generator seeded by $(0,\text{split id},i)$ and rounded to $0.1^\circ$ (split ids: train 0, val 1, test 2).

\paragraph{Neighbour tilt (sorghum).} Sorghum stems lean. To keep the three stems from closing in on each other along
the row, each neighbour is rotated about its stem base by the smallest rotation that makes its stem direction, after its
own yaw, parallel to the main plant's stem direction after the main yaw. Stem directions are the principal axis (SVD) of
the stem vertices. The main plant is never tilted; the tilt angles are recorded per orientation. Maize is not tilted.

\subsection{Contact solver}
Before rendering, the three plants are made free of intersections while staying as close as possible to their grown
shapes. Plants may touch but may not cross.

\paragraph{Leaf model.} Each leaf blade is a chain of 8\,cm segments from collar to tip. Each joint can bend its segment
along a circular arc, up/down and sideways, so the blade stays smooth, keeps its length and stays attached to the stem.
Stems and sheaths (everything within 2.5\,cm of the stem axis) never move; the maize tassel is rigid. A leaf's joint
bends add up to a total turn limit. A leaf may not bend into its own stem, and may not cross its own plant's leaves more
than it does as grown.

\paragraph{Resolution order.} Where a leaf would cross another plant, it bends at the joint before the contact by the
smallest angle that clears, trying small bends first, so it comes to rest on or under what it meets. Neighbours give way
first. A main-plant leaf bends only where a neighbour cannot get out of its way (the neighbour's stem or sheath, or a
neighbour leaf that found no clear shape), so the main plant keeps its grown, ground-truth shape wherever possible. For
sorghum, a neighbour leaf that cannot bend clear may also roll about its stem by up to $\pm90^\circ$ (candidates every
$20^\circ$, nearest first); rolled neighbour leaves are exempt from the own-plant rule. Maize leaves do not roll.

\paragraph{Exactness.} Candidate shapes are screened on 1\,cm occupancy grids; every remaining contact is then decided
by exact triangle--triangle tests on the meshes, so leaves can rest against each other with no margin. After a layout
is found it is settled: each bent leaf eases back toward its grown shape to the least bend that still clears the other
plants by exact tests.

\paragraph{Staged limits.} If a layout cannot be solved under the usual bend limits, looser limits are tried:

\begin{center}\small
\begin{tabular}{lcccc}
\toprule
Stage & collar joint sideways & blade joints sideways & up/down & total turn \\
\midrule
0 & $30^\circ$ & $10^\circ$ & $45^\circ$ & $90^\circ$ \\
1 & $60^\circ$ & $10^\circ$ & $60^\circ$ & $120^\circ$ \\
2 & $90^\circ$ & $30^\circ$ & $70^\circ$ & $180^\circ$ \\
\bottomrule
\end{tabular}
\end{center}

\paragraph{Spacing.} The solver first tries $d=0.15$\,m. If no stage solves it, the spacing is raised in 5\,cm steps up
to 1.5\,m until a stage succeeds, then refined in 1\,cm steps between the last failure and that success. The spacing
used, whether it was raised, the stage, the number of bent and rolled leaves per plant, the largest leaf turn and the
solve time are recorded per orientation.

\paragraph{Independent check.} \file{verify_cross.py} re-tests solved rows with an independent exact triangle--triangle
intersection test in NumPy on the positions as rendered; it found no crossing plant pairs on the rows checked.

\subsection{Rover camera and rendering}
One shot is taken per orientation. The rover camera sits at $(0,h,\text{dist})$ and looks horizontally at $(0,h,0)$,
where $h$ is the height of the main plant's bounding-box centre (a level camera aimed at the main stem). The camera
distance is $\text{dist}=1.0$\,m for maize and $2.5$\,m for sorghum; near $=0.05$\,m, far $=\text{dist}+2.0$\,m. The lens,
background and the three directional lights (aimed at the main plant's bounding-box centre) are those of the single-view
captures. Maize uses the maize materials and the headlight of \file{capture_maize.py}; sorghum uses the single leaf
texture of \file{capture.py} (roughness 1.0, single-sided) and no headlight.

Each shot is rendered twice from the same camera: the \emph{input} with all three plants, and the \emph{ground truth}
with the main plant alone, in its bent shape from that layout.

\subsection{Outputs}
Per shot (\file{<split>/<main>/theta_TTT/shot_000/}):

\begin{center}\small
\begin{tabularx}{\linewidth}{>{\raggedright\arraybackslash}p{4.2cm}Y}
\toprule
File & Content \\
\midrule
\file{rgba.png} & input image, all three plants; alpha 255 on plant pixels, 0 on background \\
\file{depth.png} & input depth, encoding as in the single-view sets with this shot's near/far \\
\file{pointcloud_cam.ply} & input cloud: 24{,}576 points back-projected from \file{depth.png} (Eq.~\ref{eq:backproject}), camera frame, all visible plants \\
\file{gt_rgba.png}, \file{gt_depth.png} & main plant alone, same camera \\
\file{gt_pointcloud_cam.ply} & main plant alone: 8{,}192 points, camera frame \\
\file{camera_pose.json} & position, quaternion, camera-to-world, world-to-camera, intrinsics, near, far, depth encoding, point counts, rover pose, row (plants, spacing used, angles) and bend summary \\
parameter file & main plant's parameters: maize \file{maize_XXXX_spline.xml} (the trimmed file of the single-view set); sorghum \file{Sorghum_N_spline.yml} with the point lists removed \\
\bottomrule
\end{tabularx}
\end{center}

Per row, \file{row.json} records the three plants, the rover settings and, per orientation, the angles, spacing used,
stage, bent/rolled leaves, neighbour tilt, solve time and per-shot pixel and point counts. Each dataset root holds
\file{dataset.json} with every generation setting.

\paragraph{Point counts and sampling.} The input cloud has 24{,}576 points and the ground-truth cloud 8{,}192. Points are
drawn uniformly without replacement from the back-projected plant pixels; in maize rows a shot has about
$1.9$--$5.1\times10^5$ plant pixels for the input and $0.4$--$2.5\times10^5$ for the ground truth, so both counts are
always reached. Sorghum rows are generated at these counts directly. Maize rows were first generated with 8{,}192 input
points; their input clouds were then rebuilt at 24{,}576 points from the stored \file{depth.png} and the shot's
intrinsics (\file{maize_row/resample_input_pc.py}). This rebuild matches the original back-projection: the old
8{,}192-point clouds lie within $8.4\times10^{-8}$\,m of the rebuilt full cloud. The ground-truth clouds are unchanged.

\subsection{Maize row statistics}
\begin{center}\small
\begin{tabular}{lrrrrr}
\toprule
Split & orientations & at $d=0.15$\,m & stage 0 / 1 / 2 & largest $d$ used & 99th pct.\ $d$ \\
\midrule
train & 105{,}000 & 80{,}223 (76.4\,\%) & 81{,}679 / 16{,}138 / 7{,}183 & 0.33\,m & 0.23\,m \\
val   & 22{,}500  & 17{,}622 (78.3\,\%) & 17{,}780 / 3{,}321 / 1{,}399 & 0.28\,m & 0.23\,m \\
test  & 22{,}500  & 17{,}767 (79.0\,\%) & 17{,}824 / 3{,}284 / 1{,}392 & 0.30\,m & 0.23\,m \\
\bottomrule
\end{tabular}
\end{center}
No orientation failed. On 10 sample sorghum rows (100 orientations, 2.5\,m), 85 were solved at 0.15\,m and the rest at
0.16--0.21\,m.

%======================================================================
\section{Reproduction}
\begin{verbatim}
cd /media/microway/back_up_2/RoverRows/tools
PY=/home/microway/miniconda3/envs/mvcap/bin/python

# maize rows (8,192-point input), then 24,576-point input copy into maize_row/
PYOPENGL_PLATFORM=egl $PY rover_run.py --split train --rows 0:10500 \
    --skip-done --workers 30 --gpu 0
$PY ../maize_row/resample_input_pc.py --splits train val test --workers 8

# sorghum rows
OUT=/media/microway/back_up_2/RoverRowsSorghum/Sorghum_row
PYOPENGL_PLATFORM=egl $PY rover_run.py --species sorghum --out $OUT \
    --split-dir $OUT/splits --dist 2.5 --points 24576 --points-gt 8192 \
    --split train --rows 0:10500 --skip-done --workers 20 --gpu 0
\end{verbatim}
\texttt{-{}-skip-done} resumes a run: rows whose \file{row.json} exists are skipped, and \file{row.json} is written
atomically after every shot of the row.

%======================================================================
\section{Notes and caveats}
\begin{itemize}
\item The row ground truth is the main plant \emph{as bent} in that row, not its free-standing shape. The main plant's
leaves are bent only where unavoidable; its parameter file describes the grown plant.
\item At 2.5\,m the $40^\circ$ lens covers about 1.8\,m of height, so in the sorghum single-view and row sets the top or
the base of tall plants can fall outside the frame (in the 10 sample sorghum rows the main plant touches the top edge in
58 of 100 shots).
\item The maize and sorghum single-view captures differ in view sampling (random per plant vs.\ fixed Fibonacci), camera
distance (fitted vs.\ fixed), materials and lighting (headlight in maize only), and in the files shipped (the sorghum set
has no back-projected cloud and no mesh).
\item Sorghum plant 0 uses seed 0, for which the generator does not reseed its random generator.
\todo{Check that Sorghum\_0's parameters match its mesh.}
\item The sorghum split rule is inferred from \file{assignment.csv} (see Section~3).
\end{itemize}

\end{document}
